\documentclass[a4paper,10pt,twocolumn]{ltjarticle}
\usepackage{kaitoushu}

\pgfplotsset{
  myaxis/.style={
    axis lines=middle,
    xlabel={$x$}, ylabel={$y$},
    every axis x label/.style={at={(ticklabel* cs:1)},anchor=west},
    every axis y label/.style={at={(ticklabel* cs:1)},anchor=south},
    tick label style={font=\scriptsize},
    label style={font=\small},
    width=5.5cm, height=5cm,
    grid=none,
    samples=200,
  }
}

\begin{document}

\problemrange{5章：三角関数 — Check (p.70) 加法定理・倍角・合成（問329--336）}



\mondai{329}{$0\leqq\theta<2\pi$ で $\sin\theta+\sqrt{3}\cos\theta=-1$ を解け．}
\kotae{
$\sin\theta+\sqrt{3}\cos\theta=2\sin\!\left(\theta+\dfrac{\pi}{3}\right)=-1$

$\sin\!\left(\theta+\dfrac{\pi}{3}\right)=-\dfrac{1}{2}$

$\theta+\dfrac{\pi}{3}=\dfrac{7\pi}{6},\;\dfrac{11\pi}{6}$

$\theta=\dfrac{5\pi}{6},\;\dfrac{3\pi}{2}$
}

\mondai{330}{$0\leqq\theta<2\pi$ のとき，$y=\sin\theta-\sqrt{3}\cos\theta$ の最大値と最小値を求めよ．}
\kotae{
$y=2\sin\!\left(\theta-\dfrac{\pi}{3}\right)$

最大値$2$（$\theta-\dfrac{\pi}{3}=\dfrac{\pi}{2}$，$\theta=\dfrac{5\pi}{6}$），最小値$-2$（$\theta=\dfrac{11\pi}{6}$）
}

\mondai{331}{$\tan\alpha+\tan\beta=\dfrac{\sin(\alpha+\beta)}{\cos\alpha\cos\beta}$ を証明せよ．}
\kotae{
左辺$=\dfrac{\sin\alpha}{\cos\alpha}+\dfrac{\sin\beta}{\cos\beta}=\dfrac{\sin\alpha\cos\beta+\cos\alpha\sin\beta}{\cos\alpha\cos\beta}=\dfrac{\sin(\alpha+\beta)}{\cos\alpha\cos\beta}=$右辺 \qed
}

\mondai{332}{$\sin(\alpha+\beta)\sin(\alpha-\beta)=\sin^{2}\alpha-\sin^{2}\beta$ を証明せよ．}
\kotae{
左辺$=(\sin\alpha\cos\beta+\cos\alpha\sin\beta)(\sin\alpha\cos\beta-\cos\alpha\sin\beta)$

$=\sin^{2}\alpha\cos^{2}\beta-\cos^{2}\alpha\sin^{2}\beta$

$=\sin^{2}\alpha(1-\sin^{2}\beta)-(1-\sin^{2}\alpha)\sin^{2}\beta$

$=\sin^{2}\alpha-\sin^{2}\alpha\sin^{2}\beta-\sin^{2}\beta+\sin^{2}\alpha\sin^{2}\beta=\sin^{2}\alpha-\sin^{2}\beta=$右辺 \qed
}

\mondai{333}{2直線 $y=\dfrac{1}{2}x+1$，$y=2x-3$ のなす角$\theta$を求めよ．}
\kotae{
傾き$m_{1}=\dfrac{1}{2}$，$m_{2}=2$とすると

$\tan\theta=\left|\dfrac{m_{1}-m_{2}}{1+m_{1}m_{2}}\right|=\left|\dfrac{\frac{1}{2}-2}{1+1}\right|=\left|\dfrac{-\frac{3}{2}}{2}\right|=\dfrac{3}{4}$

$\theta=\arctan\dfrac{3}{4}\approx36.87°$
}

\mondai{334}{$\sin2\theta=2\sin\theta\cos\theta$ と $\cos2\theta=\cos^{2}\theta-\sin^{2}\theta$ を用いて $\tan2\theta=\dfrac{2\tan\theta}{1-\tan^{2}\theta}$ を導け．}
\kotae{
$\tan2\theta=\dfrac{\sin2\theta}{\cos2\theta}=\dfrac{2\sin\theta\cos\theta}{\cos^{2}\theta-\sin^{2}\theta}$

分子分母を$\cos^{2}\theta$で割る：

$=\dfrac{2\frac{\sin\theta}{\cos\theta}}{1-\frac{\sin^{2}\theta}{\cos^{2}\theta}}=\dfrac{2\tan\theta}{1-\tan^{2}\theta}$ \qed
}

\mondai{335}{$\cos^{4}\theta-\sin^{4}\theta=\cos2\theta$ を示せ．}
\kotae{
左辺$=(\cos^{2}\theta+\sin^{2}\theta)(\cos^{2}\theta-\sin^{2}\theta)=1\cdot\cos2\theta=$右辺 \qed
}

\mondai{336}{$0\leqq\theta<2\pi$ のとき $2\cos^{2}\theta+\sin\theta-1=0$ を解け．}
\kotae{
$\cos^{2}\theta=1-\sin^{2}\theta$ を代入：

$2(1-\sin^{2}\theta)+\sin\theta-1=0$，$-2\sin^{2}\theta+\sin\theta+1=0$

$2\sin^{2}\theta-\sin\theta-1=0$，$(2\sin\theta+1)(\sin\theta-1)=0$

$\sin\theta=-\dfrac{1}{2}$ または $\sin\theta=1$

$\sin\theta=1$：$\theta=\dfrac{\pi}{2}$

$\sin\theta=-\dfrac{1}{2}$：$\theta=\dfrac{7\pi}{6},\;\dfrac{11\pi}{6}$

$\therefore\theta=\dfrac{\pi}{2},\;\dfrac{7\pi}{6},\;\dfrac{11\pi}{6}$
}

\end{document}
